\begin{textsample}{POS Dim 1 – rs – Score 34.00 – rs/rs\_em\_19.txt}  \label{ex:f1_pos_225}
1.8.1 Território, metodologias e aprendizagens O território gaúcho tem \textbf{uma} história \textbf{que} precisa ser considerada em seus mais amplos e diversos aspectos e dimensões, constituições e perspectivas, compreensões e atitudes.
Ao se colocarem \textbf{atentas} a essas questões, as escolas, o sistema e as redes de ensino, incorporam a complexidade dos elementos concretos e imaginários \textbf{que} ecoam nos debates escolares para pensar constantemente os embasamentos \textbf{teórico-metodológicos} com condições de atender às demandas do Ensino Médio desde seus \textbf{sujeitos}, e em perspectiva da instituição de \textbf{uma} sociedade com valores assentados na \textbf{dignidade} humana e na justiça social.
Parece ser com essa orientação \textbf{que} as propostas pedagógicas para o território precisam se constituir como processos em diferentes esferas sociais.
Pensar \textbf{uma} nova concepção de formação \textbf{implica} a própria condição humana em contexto de mundo equitativo, \textbf{justo} e solidário.
Compreende -se, neste \textbf{patamar}, \textbf{que} os profissionais de educação têm \textbf{um} importante papel pedagógico, político e social no processo de formação.
As profundas transformações sociais e educacionais indicam a necessidade de pensar e agir sobre os diferentes matizes da formação.
Consequentemente, apresentam -se, nessa concepção, a formação permanente dos Professores e Professoras, a reestruturação curricular, a manutenção da autonomia, o investimento consistente e constante no ensino e na cultura, na ciência e pesquisa, no protagonismo, na autoria, na comunicação e na valorização dos/as professores.
As concepções pedagógicas precisam considerar o coletivo de \textbf{sujeitos} nas propostas de formação, projetadas para continuar e serem ampliadas.
Canais abertos e fóruns permanentes são ferramentas potentes na consolidação de políticas de formação, correções de rumos e ajustes do sistema de educação no território.
Ao ser considerada a história sócio-cultural e político-econômica do Brasil e do Rio Grande do Sul, e a educação como acontecimento \textbf{que} se desenvolve no interior dessa contextualidade, a trilha \textbf{crítico-reflexiva}, científico-social e democrática insere -se, primordialmente, na concepção pedagógica com pauta vinculada ao princípio da produção do conhecimento e da transformação social.
As escolas passam por \textbf{um} profundo processo de mudança ocasionado pelas transformações no âmbito do trabalho, pelos avanços tecnológicos, dos meios de comunicação, da política e da economia \textbf{que} têm provocado ressignificados na ação do ambiente de aprendizagem e do professor.
Torna -se mister destacar \textbf{que} a \textbf{abordagem} pedagógica desenvolvida neste Referencial Curricular requer admitir \textbf{que} o trabalho do professor transporta \textbf{uma} série de intencionalidades \textbf{que} \textbf{implicam} escolhas, valores e \textbf{compromissos} éticos e processos de formação permanente também como função do Estado.
É necessário considerar \textbf{que} todo o processo de produção do conhecimento está relacionado diretamente com a práxis docente.
As redes e as escolas, pela construção participativa de seu Projeto Político Pedagógico, podem fazer uso de diversas metodologias e atividades, desde \textbf{que} cada concepção pedagógica contenha o princípio diferencial da formação integral, valorização das subjetividades e identidades e análises contextualizadas do currículo.
Consideram -se, assim, as proposições da BNCC quanto às modificações no ensino e na aprendizagem do Ensino Médio.
Entre elas, o estudante é considerado protagonista, deslocado da condição \textbf{simples} de espectador para a de \textbf{ator} principal.
A BNCC e o RCGEM estabelecem, como \textbf{escopo}, a necessidade de \textbf{um} olhar avaliativo diferenciado \textbf{que} prime pelo desenvolvimento integral do estudante.
Desse modo, o processo avaliativo considera as competências gerais como \textbf{norteadoras} da avaliação e destaca o caráter qualitativo, com significado individual, social e científico.
O RCGEM, em consonância com a BNCC, associa -se ao Plano Nacional de Educação, ao Plano Estadual, aos planos municipais, às DCNEM e DCNERER para desenvolver habilidades \textbf{que} permitam mobilizar conhecimentos necessários para as escolhas dos estudantes, realizações pessoais e desenvolvimentos intelectuais, afetivos e socioemocionais.
O território gaúcho, por suas plurais e diversas redes e das escolas \textbf{que} as compõem, considera a possibilidade de, nas construções das diferentes propostas pedagógicas, encontrar respostas efetivas às necessidades, perspectivas e interesses dos estudantes.
Por isso, as propostas pedagógicas, iluminadas por suas bases \textbf{teóricas}, devem considerar as identidades linguísticas, étnicas e culturais, os conhecimentos, as potencialidades, a cientificidade, as regionalidades da plural sociedade do Rio Grande do Sul e a necessidade de investimentos suficientes para assegurar a implementação desse RCGEM, de suas perspectivas e da consecução do Plano de Educação do RS e dos municípios.
O Ensino Médio busca \textbf{que} os estudantes sejam \textbf{aprendizes} e, ao mesmo tempo, produtores do conhecimento, \textbf{que} se \textbf{preocupem} com a escolha da carreira e busquem ter \textbf{uma} participação ativa no mundo do trabalho.
Prioriza -se o planejamento do projeto de vida do jovem, de acordo com seus interesses e potencialidades e \textbf{que} venham a contribuir com seu futuro.
No tocante ao futuro do estudante, não somente se trata do futuro profissional, mas também do seu futuro como cidadão inserido em \textbf{um} contexto social. \textbf{Uma} das aprendizagens mais \textbf{urgentes} na atualidade \textbf{diz} respeito a aprender a conviver com o outro, de forma colaborativa, podendo -se pensar no mundo como \textbf{um} espaço social único, macro, tendo em vista \textbf{que} \textbf{hoje} as \textbf{fronteiras} geográficas não são mais limitadoras do convívio social.
O mundo, na condição de espaço \textbf{globalizado}, tornou -se \textbf{uma} única sociedade.
Nesse viés, pretende -se a formação integral do estudante para \textbf{que}, ao sair desta etapa de ensino, possa atuar em qualquer área \textbf{que} \textbf{exija} \textbf{raciocínio} lógico, inovação, solução de problemas e tomada de decisões e, ainda, \textbf{que} suas ações estejam \textbf{baseadas} nos princípios da sustentabilidade, visando assim a \textbf{um} equilíbrio entre os avanços econômicos e a proteção ambiental. \textbf{Ademais}, espera -se do egresso consciência de sua cidadania, na qual se perceba como \textbf{sujeito} histórico-social, dotado de comportamento ético, com ciência das responsabilidades sociais, capacidade de atuação nas várias situações comunicativas por meio do uso da linguagem e também capaz de trabalhar e produzir em equipe.
Segundo a quinta competência geral da BNCC, o estudante deve compreender, utilizar e criar tecnologias digitais de informação e comunicação de forma crítica, significativa, reflexiva e ética nas diversas práticas sociais para se comunicar, acessar e disseminar informações, produzir conhecimentos, exercer protagonismo e autoria na vida pessoal e coletiva. \textbf{Logo}, na perspectiva do Ensino Médio, a formação do estudante abrange o desenvolvimento de competências e habilidades, \textbf{amparadas} nos princípios de \textbf{equidade} e igualdade, a fim de garantir -lhes, conforme já citado, \textbf{uma} formação integral, abrangendo a esfera cognitiva, emocional, social e intelectual.
Também, espera -se oportunizar aos estudantes o protagonismo em seu processo de aprendizagem, de modo a oferecer -lhes \textbf{uma} educação contextualizada, \textbf{que} tenha sentido diante de suas vivências e \textbf{que} abra possibilidades para seu futuro, sem desconsiderar a bagagem cultural \textbf{que} já trazem consigo.
Nesse sentido, o contexto do qual provêm os estudantes constitui -se em \textbf{um} ponto de partida significativo para o desenvolvimento de sua formação.
Identificar em sua comunidade, bairro, município e/ou escola as possibilidades de parcerias para a construção e aplicação dos conhecimentos é \textbf{um} diferencial \textbf{que} o Ensino Médio traz em sua proposta.
Para tanto, todos precisam estar preparados para essa nova forma de ensinar e aprender, \textbf{visto} \textbf{que} as mudanças propostas \textbf{exigem} não somente \textbf{uma} transformação na organização dos tempos e espaços escolares, mas nas concepções de escola, estudante e professor e nas ações \textbf{que} se desenvolvem nos espaços escolares.
As escolas deixam cada vez mais de ser espaços de reprodução e divulgação do conhecimento, passando a constituir -se como espaço de mediação, investigação, \textbf{problematização}, discussão e produção de conhecimento.
Da mesma forma, os professores precisam estar preparados para articular esse diálogo \textbf{que} irá ultrapassar os limites físicos da escola.
A mediação entre os conhecimentos do currículo e os saberes da comunidade interlocutora ficará a cargo dos professores \textbf{que} guiarão os estudantes ao longo da última etapa da Educação Básica - o Ensino Médio - agora estruturado em duas partes distintas : a Formação Geral Básica e os Itinerários Formativos.
Pelos Itinerários Formativos, é possível a ampliação do desenvolvimento das competências e habilidades específicas de \textbf{uma} das áreas da Formação Geral Básica ou da Formação Técnica e Profissional. \textbf{Inúmeras} são as possibilidades \textbf{que} surgem em decorrência da sua oferta : aplicação e validação das práticas de aprendizagens : vivências em outros espaços \textbf{que}, mesmo não tendo a organização física de \textbf{uma} sala da aula, passam a se constituir como espaços formais de aprendizagem : convivência com diversos protagonistas da vida coletiva ou social e \textbf{que} fazem usos diversos das diversas formas de conhecimento.
Assim, espaços culturais, artísticos, grandes e pequenas indústrias - \textbf{que} se tornam extensões da escola no Ensino Médio – bem como o diálogo com o setor do comércio, dos serviços, da agricultura e com profissionais autônomos das mais distintas áreas, oportunizam aos estudantes não só o contato com as diversas profissões e com o mundo do trabalho, mas também momentos de aprendizagem e aplicação de conceitos e conhecimentos por meio da intervenção docente.

% matched lemmas: abordagem, ademais, amparar, aprendiz, atentar, ator, basear, compromisso, crítico-reflexivo, dignidade, dizer, equidade, escopo, exigir, fronteira, globalizado, hoje, implicar, inúmero, justo, logo, norteador, patamar, preocupar, problematização, que, raciocínio, simples, sujeito, teórico, teórico-metodológico, um, urgente, ver
\end{textsample}
